\subsection{Earlier Point-Count Controls and Protocol Dependence}
\label{sec:mn-history}

The completed work extends beyond the final twelve-branch comparison.
The HPC inventory records eight earlier classifier jobs in addition to
those branches and the two mesh-reference controls. Table~\ref{tab:mn-history}
retains their actual model input sizes; a 512-point file resampled to
1,024 rows is distinct from a native 512-point input.

\begin{table}[!t]\centering\setlength{\tabcolsep}{4pt}
\begin{tabular}{L{0.46\linewidth}rrr}\toprule
Earlier representation & Input points & Best OA & Final OA\\\midrule
Original baseline & 1,024 & 91.95 & ---\\
Sparse 50\%, native input & 512 & 91.26 & 91.14\\
Sparse 50\%, loader resampled & 1,024 & 91.17 & ---\\
Sparse 50\%, EAR $\times2$ & 1,024 & 90.82 & ---\\
Sparse 50\%, EAR $\times4$ & 2,048 & 90.61 & 89.76\\
Sparse 50\%, PDANS $\times4$ & 2,048 & 91.47 & 90.99\\
Original, EAR $\times4$ & 4,096 & 91.48 & 90.85\\
Original, PDANS $\times4$ & 4,096 & 91.62 & 90.87\\\bottomrule
\end{tabular}
\caption{Eight completed earlier HPC PointNet++ jobs. OA is the recorded
batch aggregate in percent. Dashes indicate final-epoch values absent
from this inventory, not failed evaluations. Earlier/final repetitions
of the original, EAR 4,096 and PDANS 4,096 branches have identical
79/79 weight tensors at seed 42 and are not independent training seeds.}\label{tab:mn-history}
\end{table}

PDANS illustrates why the earlier results cannot be silently merged with
the final sparse-input matrix. At 512 visible points, its 2,048-point
output achieved 91.47\% Best OA against a native baseline of 91.26\%,
a descriptive gain of 0.21~pp. At 256 visible points, its 1,024-point
output achieved 90.34\% against 90.85\%, a loss of 0.51~pp. The ratio
remained four, but both the visible support and downstream input size
changed. This is a protocol-dependent sign change, not a controlled
estimate of sparsity alone or proof of a statistically reliable gain.

The 1,024-row controls further separate nominal input size from acquired
information. The genuine original representation reached 91.95\%,
the loader-resampled 512-point representation 91.17\%, and EAR's
512-to-1,024 output 90.82\%. The 0.35~pp advantage of loader resampling
over this EAR-style output is limited to these trained runs. It does
not imply that duplicating points universally improves classification.
Together with the native-input and mesh-reference controls, these jobs
explain the transition to a fixed fourfold, two-line final protocol.
